\section{Metodología}

El desarrollo experimental de la práctica se estructuró en cuatro etapas secuenciales, cumpliendo estrictamente con las normas de seguridad del laboratorio y las pautas técnicas recomendadas.

\subsection{Montaje Inicial y Medición en Frío}
\begin{enumerate}
    \item Inspeccionar física y eléctricamente los bornes y aislamientos del transformador monofásico de prueba, verificando los datos de placa y las conexiones de tierra de seguridad.
    \item Registrar la temperatura ambiente inicial $T_1$ con el termómetro de precisión ubicado en las inmediaciones del puesto de trabajo.
    \item Con el equipo totalmente desenergizado y habiendo permanecido en reposo térmico suficiente para que la masa conductora se encuentre a la temperatura ambiente, conectar el instrumento de medición de resistencia en corriente continua a los bornes secundarios ($X_1 - X_2$). Medir y registrar la resistencia en frío $R_2(T_1)$ (o $R_f$).
    \item Realizar el alambrado del montaje principal según el esquema de la Figura~\ref{fig:conexion_ensayo_carga}, verificando que los breakers de protección permanezcan abiertos y el VARIAC esté posicionado en su punto de tensión cero.
    \item Solicitar la inspección y el visto bueno del profesor o técnico de laboratorio antes de proceder a la energización del circuito.
\end{enumerate}

\subsection{Parte A: Ensayo de Rendimiento con Carga Resistiva Pura (\texorpdfstring{$\cos\varphi = 1$}{cos phi = 1})}
\begin{enumerate}
    \item Cerrar los interruptores de protección y alimentar el transformador por su devanado de baja tensión (primario en configuración elevadora) mediante el VARIAC hasta alcanzar de manera precisa la tensión nominal $V_1 = V_{1n} = 110\text{ V}$.
    \item Con la carga en circuito abierto ($I_2 = 0\text{ A}$), registrar los valores de vacío: $V_1, I_1, P_1, V_2, I_2, P_2$. Este punto permite validar las pérdidas en el núcleo $P_{\text{Fe}}$ determinadas previamente en el ensayo de vacío de la Práctica Nº 4.
    \item Conectar los instrumentos de medición de potencia: para el primario, conectar el vatímetro electrodinámico a sus terminales de $5\text{ A}$ y $120\text{ V}$, incorporando un transformador de corriente de relación $2:1$ para la bobina amperimétrica (protegiendo el instrumento y adaptando el rango de corriente hasta $10\text{ A}$); en el secundario, conectar el vatímetro a sus terminales de $5\text{ A}$ y $240\text{ V}$.
    \item Conectar el banco de cargas resistivas en el lado de alta tensión (secundario) y variar gradualmente la resistencia en pasos sucesivos correspondientes a los diferentes niveles de potencia ensayados.
    \item En cada paso de carga, reajustar con el VARIAC la tensión primaria de entrada para garantizar que se mantenga en el entorno de los $110\text{ V}$, compensando las caídas de tensión en las líneas de alimentación.
    \item Esperar la estabilización de los instrumentos y registrar simultáneamente las lecturas de los seis instrumentos: tensión primaria ($V_1$), corriente primaria ($I_1$), potencia activa de entrada ($P_1$), tensión secundaria ($V_2$), corriente de carga ($I_2$) y potencia activa de salida ($P_2$).
\end{enumerate}

\subsection{Parte B: Medición de Temperatura por el Método de la Resistencia}
\begin{enumerate}
    \item Mantener el transformador operando a plena carga durante el tiempo estipulado para propiciar el calentamiento representativo de los devanados.
    \item Cumplido el régimen de plena carga, desenergizar de inmediato el circuito abriendo el interruptor termomagnético de alimentación y desconectar rápidamente el banco de cargas secundario.
    \item Conectar sin demora el instrumento de corriente continua a los bornes del secundario y registrar el valor de la resistencia en caliente $R_c = R_{2c}$, minimizando el tiempo transcurrido para evitar el enfriamiento térmico del arrollamiento (siguiendo las recomendaciones del punto 5.7 de la norma COVENIN 3172-95).
    \item Registrar nuevamente la temperatura ambiente del recinto para obtener el promedio de la prueba.
    \item Calcular la elevación de temperatura y la temperatura final del bobinado $T_2$ mediante la ley de extrapolación lineal del cobre (ecuación~\eqref{eq:metodo_resistencia}).
    \item Corregir la resistencia equivalente $R_{\text{eq2}}$ a la temperatura de régimen $T_2$ (ecuación~\eqref{eq:correccion_resistencia}) y reevaluar el rendimiento a plena carga.
\end{enumerate}

\subsection{Parte C: Ensayo de Rendimiento con Carga Inductiva (\texorpdfstring{$\cos\varphi < 1$}{cos phi < 1})}
\begin{enumerate}
    \item Conectar en serie con la resistencia variable la bobina de reactancia inductiva calibrada para fijar un factor de potencia atrasado ($\cos\varphi \approx 0{,}8$).
    \item Reanudar la alimentación ajustando el VARIAC para mantener $V_1 = V_{1n}$.
    \item Ajustar la impedancia de carga hasta situar la corriente secundaria en las proximidades del punto de rendimiento máximo calculado ($I_2 \approx I_2^* = \sqrt{P_{\text{Fe}}/R_{\text{eq2}}}$).
    \item Registrar los valores medidos en los seis instrumentos de prueba ($V_1, I_1, P_1, V_2, I_2, P_2$).
    \item Determinar experimentalmente el factor de potencia secundario mediante la relación $\cos\varphi = P_2 / (V_2 I_2)$.
    \item Desenergizar el sistema y efectuar la medición de resistencia en corriente continua para corroborar la temperatura térmica en este régimen.
\end{enumerate}
